\begin{proof}[Proof of \cref{obj:prop:coincident-policy-reduction}]
We use epochs \(1,\ldots,T\), shifting the zero-based epochs of \cref{obj:lem:observable-intervention-moments} by one. Every sum over an empty range is understood to be zero.

\begin{enumerate}
\item Since \(e=b\), substitution in the policy-induced transition matrices gives, for every \(s=(x,h)\in\mathcal S\) and \(s'\in\mathcal S\),
\[
P_e(s,s')
=
\sum_{a\in\{0,1\}}e(a\mid x)
   \int_{[0,1]}K(dy,s'\mid s,a)
=
\sum_{a\in\{0,1\}}b(a\mid x)
   \int_{[0,1]}K(dy,s'\mid s,a)
=
P_b(s,s').
\]
The stationary-law assignment evaluated at these identical matrices likewise gives
\[
d_e=d_b.
\]
% lean: CausalSmith.Stat.PomdpBinaryhiddenRate.coincidentPolicy_kernel_stationary

\item We first verify membership in the class of \cref{obj:def:binary-pomdp-class} at logarithmic overlap parameter \(\zeta=0\). Sequential assignment and \(e=b\) give
\[
e(a\mid x)=b(a\mid x)\geq0,
\qquad
\sum_{a\in\{0,1\}}e(a\mid x)=1,
\qquad
e(a\mid x)\leq \exp(0)b(a\mid x)=b(a\mid x).
\]
Thus \cref{obj:ass:policy-overlap} holds with \(L=1\). Moreover, for every \(\mu,\mu'\in\Delta(\mathcal S)\),
\[
\|\mu P_e-\mu'P_e\|_{\mathrm{TV}}
=
\|\mu P_b-\mu'P_b\|_{\mathrm{TV}}
\leq
\alpha\|\mu-\mu'\|_{\mathrm{TV}},
\]
so \cref{obj:ass:target-contraction} follows from \cref{obj:ass:behavior-contraction}. Together with the remaining hypotheses, these facts establish membership in \(\mathcal M_T^{(2)}(t_0,0)\).

Write the target-policy conditional mean reward vector as
\[
r_e(s):=
\sum_{a\in\{0,1\}}e(a\mid x)
\sum_{s'\in\mathcal S}\int_{[0,1]}y\,K(dy,s'\mid s,a),
\qquad s=(x,h)\in\mathcal S.
\]
Applying \cref{obj:lem:observable-intervention-moments} at depth zero gives, at every observed epoch,
\[
m_0=d_bP_e^0r_e=d_br_e=d_er_e=\theta,
\]
where the last equality is \cref{obj:def:target-value}.
% lean: CausalSmith.Stat.PomdpBinaryhiddenRate.observable_intervention_moment_zero

\item To identify this moment with the ordinary reward mean, define the observed action likelihood ratio and depth-zero score by
\[
\rho_t:=
\begin{cases}
e(A_t\mid X_t)/b(A_t\mid X_t),& b(A_t\mid X_t)>0,\\
0,& b(A_t\mid X_t)=0,
\end{cases}
\qquad
Z_t(0):=Y_t\rho_t,
\qquad 1\leq t\leq T.
\]
Here probability and expectation are taken under the behavior experiment. By \cref{obj:ass:sequential-ignorability},
\[
\Pr_b\!\left(b(A_t\mid X_t)=0\right)
=
\mathbb E_b\!\left[
\sum_{\substack{a\in\{0,1\}\\ b(a\mid X_t)=0}}
b(a\mid X_t)
\right]
=0.
\]
On the complementary event, \(e=b\) gives \(\rho_t=1\). Consequently,
\[
Z_t(0)=Y_t\quad\text{almost surely},
\qquad
\mathbb E_b[Y_t]
=
\mathbb E_b[Z_t(0)]
=
m_0
=
\theta.
\]
% lean: CausalSmith.Stat.PomdpBinaryhiddenRate.coincidentPolicy_reward_mean

\item Define the sample mean by
\[
\overline Y_T:=T^{-1}\sum_{t=1}^{T}Y_t.
\]
The bounded rewards are integrable, so linearity of expectation and \(T\geq1\) yield
\[
\mathbb E_b[\overline Y_T]
=
T^{-1}\sum_{t=1}^{T}\mathbb E_b[Y_t]
=
T^{-1}T\theta
=
\theta.
\]
% lean: CausalSmith.Stat.PomdpBinaryhiddenRate.coincidentPolicy_sampleMean_unbiased

\item The reward support and \cref{obj:ass:pomdp-kernel} give
\[
\Pr_b(Y_t\in[0,1])
=
\mathbb E_b\!\left[
K([0,1]\times\mathcal S\mid S_t,A_t)
\right]
=1.
\]
Thus every reward is square integrable. For square-integrable real random variables \(W_1,W_2\) under the behavior law, use
\[
\operatorname{Cov}_b(W_1,W_2)
:=
\mathbb E_b\!\left[
(W_1-\mathbb E_b[W_1])(W_2-\mathbb E_b[W_2])
\right],
\qquad
\operatorname{Var}_b(W_1)
:=
\operatorname{Cov}_b(W_1,W_1).
\]
Since \(t_0>0\),
\[
0\leq\alpha=\exp(-1/t_0)<1.
\]
Also \(Y_t^2\leq1\) almost surely, and hence
\[
\mathbb E_b[Y_t^2]\leq\mathbb E_b[1]=1,
\qquad
\operatorname{Cov}_b(Y_t,Y_t)
=
\mathbb E_b[Y_t^2]-(\mathbb E_b[Y_t])^2
\leq1.
\]
% lean: CausalSmith.Stat.PomdpBinaryhiddenRate.coincidentPolicy_sampleMean_mse

\item Expanding the variance of the reward sum gives
\[
\operatorname{Var}_b\!\left(\sum_{t=1}^{T}Y_t\right)
=
\sum_{t=1}^{T}\sum_{u=1}^{T}
\operatorname{Cov}_b(Y_t,Y_u).
\]
For \(1\leq t<u\leq T\), \cref{obj:lem:observable-intervention-moments}, applied at depth zero in the class established above, supplies
\[
\left|
\mathbb E_b\!\left[
(Z_t(0)-m_0)(Z_u(0)-m_0)
\right]
\right|
\leq\alpha^{u-t-1}.
\]
Using the almost-sure score identities and reward means already established,
\[
\operatorname{Cov}_b(Y_t,Y_u)
=
\mathbb E_b\!\left[
(Z_t(0)-m_0)(Z_u(0)-m_0)
\right],
\]
and therefore
\[
\operatorname{Cov}_b(Y_t,Y_u)
\leq
\left|\operatorname{Cov}_b(Y_t,Y_u)\right|
\leq\alpha^{u-t-1}.
\]
% lean: CausalSmith.Stat.PomdpBinaryhiddenRate.coincidentPolicy_reward_covariance

To bound each covariance row, first use the finite geometric identity
\[
(1-\alpha)\sum_{\ell=0}^{T-1}\alpha^\ell
=
1-\alpha^T
\leq1,
\qquad
\sum_{\ell=0}^{T-1}\alpha^\ell
\leq\frac1{1-\alpha}.
\]
Here \(\ell\) is a nonnegative integer exponent. For fixed \(t\), the exponent maps \(u\mapsto t-u-1\) on \(1\leq u<t\) and \(u\mapsto u-t-1\) on \(t<u\leq T\) are injective and take values in \(\{0,\ldots,T-1\}\). Because all powers of \(\alpha\) are nonnegative, symmetry of covariance gives
\[
\begin{aligned}
\sum_{u=1}^{t-1}\operatorname{Cov}_b(Y_t,Y_u)
&=
\sum_{u=1}^{t-1}\operatorname{Cov}_b(Y_u,Y_t)
\leq
\sum_{u=1}^{t-1}\alpha^{t-u-1}
\leq
\sum_{\ell=0}^{T-1}\alpha^\ell
\leq\frac1{1-\alpha},\\
\sum_{u=t+1}^{T}\operatorname{Cov}_b(Y_t,Y_u)
&\leq
\sum_{u=t+1}^{T}\alpha^{u-t-1}
\leq
\sum_{\ell=0}^{T-1}\alpha^\ell
\leq\frac1{1-\alpha}.
\end{aligned}
\]
These inequalities include the empty past or future ranges at the endpoints. Splitting the row into its diagonal, past, and future parts now yields
\[
\begin{aligned}
\sum_{u=1}^{T}\operatorname{Cov}_b(Y_t,Y_u)
&=
\operatorname{Cov}_b(Y_t,Y_t)
+\sum_{u=1}^{t-1}\operatorname{Cov}_b(Y_t,Y_u)
+\sum_{u=t+1}^{T}\operatorname{Cov}_b(Y_t,Y_u)\\
&\leq1+\frac2{1-\alpha}.
\end{aligned}
\]
Summing over \(t\) proves
\[
\operatorname{Var}_b\!\left(\sum_{t=1}^{T}Y_t\right)
\leq
T\left(1+\frac2{1-\alpha}\right).
\]
% lean: CausalSmith.Stat.PomdpBinaryhiddenRate.geometricCovariance_row_sum_le

\item Finally, unbiasedness identifies the mean squared error with the variance. Rescaling the preceding bound, using \(T>0\), gives
\[
\begin{aligned}
\mathbb E_b\!\left[(\overline Y_T-\theta)^2\right]
&=
\operatorname{Var}_b(\overline Y_T)\\
&=
T^{-2}\operatorname{Var}_b\!\left(\sum_{t=1}^{T}Y_t\right)\\
&\leq
T^{-2}T\left(1+\frac2{1-\alpha}\right)
=
\frac{1+2/(1-\alpha)}{T}.
\end{aligned}
\]
% lean: CausalSmith.Stat.PomdpBinaryhiddenRate.coincidentPolicy_sampleMean_mse
\end{enumerate}
\end{proof}
